% !TEX root = ../../../main.tex
% 所属：第四部分：Fourier 分析 / 第4章：$\R$上的Fourier变换
% 节：4.1 基本定义
% 正文按原文件顺序保存；由 ../chapter4.tex 调用。
\section{基本定义}

在 Fourier 分析中，我们采用 Cauchy 主值积分作为实轴上积分的一种延拓，即
\[
\lim_{N\to\infty}\int_{-N}^{N}f(x)\dd x
=(\mathrm{CPV})\int_{-\infty}^{\infty}f(x)\dd x,
\]
并将其简记为
\[
\int_{-\infty}^{\infty}f(x)\dd x.
\]
需要注意的是，这一节中的积分与 Part1 中所讲的积分并不完全相同。为了确保上述积分存在，我们对函数作如下限制。

\begin{definition}
	如果对于定义在 $\R$ 上的\textbf{连续函数} $f$，存在常数 $A>0$，使得对任意 $x\in\R$ 都有
	\[
	\abs{f(x)}\leq \frac{A}{1+x^2},
	\]
	则称 $f$ 是\textbf{适度下降的}。
	
	记 $\R$ 上全体适度下降的函数为 $\mathcal M(\R)$。在通常的函数加法与数乘下，$\mathcal M(\R)$ 是 $\C$ 上的一个向量空间。
	
\end{definition}
该定义说明 $f$ 是有界的，并且由于$\dfrac{A}{1+x^2}\leq \dfrac{A}{x^2}(x\neq0),$ $f$ 在无穷远处至少与 $\dfrac{1}{x^2}$ 具有相同的下降速度。


对于 $f\in\mathcal M(\R)$，定义
\[
\int_{-\infty}^{\infty}f(x)\dd x
:=\lim_{N\to\infty}\int_{-N}^{N}f(x)\dd x.
\]
显然，
\[
\abs{\int_{-\infty}^{\infty}f(x)\dd x}
\leq \int_{-\infty}^{\infty}\abs{f(x)}\dd x
\leq \int_{-\infty}^{\infty}\frac{A}{1+x^2}\dd x
=A\pi,
\]
所以该积分存在。

\textbf{注.}\quad 适度下降的条件也可写成：存在 $A>0$ 与 $\eps>0$，使得
\[
\abs{f(x)}\leq \frac{A}{1+\abs{x}^{1+\eps}}.
\]
为方便起见，这里取 $\eps=1$。

\begin{proposition}\label{prop:moderate-integral-properties}
	适度下降函数的积分满足如下性质：
	\begin{enumerate}[label=(\roman*)]
		\item \textbf{线性性：}若 $f,g\in\mathcal M(\R)$ 且 $a,b\in\C$，则
		\[
		\int_{-\infty}^{\infty}\bigl(af(x)+bg(x)\bigr)\dd x
		=a\int_{-\infty}^{\infty}f(x)\dd x
		+b\int_{-\infty}^{\infty}g(x)\dd x.
		\]
		
		\item \textbf{平移不变性：}对任意 $h\in\R$，有
		\[
		\int_{-\infty}^{\infty}f(x-h)\dd x
		=\int_{-\infty}^{\infty}f(x)\dd x.
		\]
		
		\item \textbf{伸缩性：}若 $\delta>0$，则
		\[
		\delta\int_{-\infty}^{\infty}f(\delta x)\dd x
		=\int_{-\infty}^{\infty}f(x)\dd x.
		\]
		
		\item \textbf{连续性：}若 $f\in\mathcal M(\R)$，则当 $h\to0$ 时，
		\[
		\int_{-\infty}^{\infty}\abs{f(x-h)-f(x)}\dd x\longrightarrow0.
		\]
	\end{enumerate}
\end{proposition}

\begin{proof}
	(i) 由积分的线性性立即得到。
	
	(ii) 不妨设 $h\geq0$。由变量替换，
	\[
	\int_{-N}^{N}f(x-h)\dd x
	=\int_{-N-h}^{N-h}f(x)\dd x.
	\]
	当 $N>h$ 时，
	\[
	\begin{aligned}
		&\abs{\int_{-N}^{N}f(x-h)\dd x-\int_{-N}^{N}f(x)\dd x}\\
		&\qquad\leq
		\int_{-N-h}^{-N}\abs{f(x)}\dd x
		+\int_{N-h}^{N}\abs{f(x)}\dd x
		\leq \frac{2Ah}{1+(N-h)^2}\longrightarrow0.
	\end{aligned}
	\]
	令 $N\to\infty$ 即得结论；$h<0$ 的情形同理。
	
	(iii) 由变量替换，
	\[
	\delta\int_{-N}^{N}f(\delta x)\dd x
	=\int_{-\delta N}^{\delta N}f(u)\dd u.
	\]
	令 $N\to\infty$ 类似于(ii)即得证。
	
	(iv) 只需考虑 $\abs{h}\leq1$。给定 $\eps>0$，取充分大的 $N$，使得
	\[
	\int_{\abs{x}\geq N-1}\abs{f(x)}\dd x<\frac{\eps}{4}.
	\]
	于是对 $\abs{h}\leq1$，
	\[
	\int_{\abs{x}\geq N}\abs{f(x-h)}\dd x
	\leq\int_{\abs{u}\geq N-1}\abs{f(u)}\dd u
	<\frac{\eps}{4},
	\]
	并且
	\[
	\int_{\abs{x}\geq N}\abs{f(x)}\dd x<\frac{\eps}{4}.
	\]
	固定 $N$。由于 $f$ 在 $[-N-1,N+1]$ 上一致连续，当 $h\to0$ 时，
	\[
	\sup_{\abs{x}\leq N}\abs{f(x-h)-f(x)}\longrightarrow0.
	\]
	因此可取充分小的 $h$，使得上述上确界小于 $\eps/(4N)$。此时
	\[
	\begin{aligned}
		\int_{-\infty}^{\infty}\abs{f(x-h)-f(x)}\dd x
		\leq{}&\int_{-N}^{N}\abs{f(x-h)-f(x)}\dd x\\
		&+\int_{\abs{x}\geq N}\abs{f(x-h)}\dd x
		+\int_{\abs{x}\geq N}\abs{f(x)}\dd x\\
		<&\frac{\eps}{2}+\frac{\eps}{4}+\frac{\eps}{4}=\eps.
	\end{aligned}
	\]
	故结论成立。
\end{proof}
另一个可能会用到概念是支集
\begin{definition}
	函数 $f$ 的支集记作
	
	$$
	\operatorname{supp} f
	$$
	
	定义为
	
	$$
	\operatorname{supp} f=\overline{\{x: f(x) \neq 0\}} .
	$$
	
	也就是：先取出所有使 $f(x) \neq 0$ 的点，再把这个集合闭包起来。
\end{definition}
\begin{definition}
	若 $f\in\mathcal M(\R)$，则在 $\xi\in\R$ 处定义 $f$ 的 Fourier 变换为
	\[
	\widehat f(\xi)
	=\int_{-\infty}^{\infty}f(x)e^{-2\pi i x\xi}\dd x.
	\]
	
\end{definition}
由于 $\abs{e^{-2\pi i x\xi}}=1$，这是一个关于适度下降函数的积分，因而积分有意义。
\begin{proposition}
	若 $f\in\mathcal M(\R)$，则 $\widehat f$ 是一个连续且有界的函数，并且当 $\abs{\xi}\to\infty$ 时，
	\[
	\widehat f(\xi)\longrightarrow0.
	\]
\end{proposition}

\begin{proof}
	首先，
	\[
	\abs{\widehat f(\xi)}
	\leq\int_{-\infty}^{\infty}\abs{f(x)}\dd x
	\leq A\pi,
	\]
	故 $\widehat f$ 有界。
	
	若 $\xi_n\to\xi$，则
	\[
	f(x)e^{-2\pi i x\xi_n}\longrightarrow f(x)e^{-2\pi i x\xi},
	\qquad
	\abs{f(x)e^{-2\pi i x\xi_n}}\leq\abs{f(x)}.
	\]
	由含参变量积分的连续性定理，$\widehat f(\xi_n)\to\widehat f(\xi)$，所以 $\widehat f$ 连续。
	
	最后，设 $\xi\neq0$，令 $h=1/(2\xi)$。由平移不变性与变量替换，
	\[
	\int_{-\infty}^{\infty}f(x-h)e^{-2\pi i x\xi}\dd x
	=e^{-2\pi i h\xi}\widehat f(\xi)
	=-\widehat f(\xi).
	\]
	因此
	\[
	\widehat f(\xi)
	=\frac12\int_{-\infty}^{\infty}
	\left(f(x)-f\left(x-\frac1{2\xi}\right)\right)e^{-2\pi i x\xi}\dd x.
	\]
	从而
	\[
	\abs{\widehat f(\xi)}
	\leq\frac12\int_{-\infty}^{\infty}
	\abs{f(x)-f\left(x-\frac1{2\xi}\right)}\dd x.
	\]
	当 $\abs{\xi}\to\infty$ 时，$1/(2\xi)\to0$，由命题 \ref{prop:moderate-integral-properties} 的连续性结论，右端趋于 $0$，故
	\[
	\widehat f(\xi)\longrightarrow0.
	\]
\end{proof}
